% SCRATCH DRAFT — Vol 2 courtroom coda, Ch16 "Matter Waves and Scattering" (task vol2-coda, author: Podo).
% THIS IS NOT THE BOOK. Do not \input. Floor commit not yet banked; 16.tex untouched.
% ⚠ 2nd PROVED LOOP RESIDUE. Fresh-append (no old beat; body ends line 206).
%
% ⚠ FENCES:
%  - 2nd PROVED LOOP RESIDUE: ANNOUNCE proved (as Ch5 — a finite theorem; at the peak the loop is charged, off
%    it trivial; nothing between; no continuum) but WITHHOLD the ±1 VALUE entirely — no number, no sign committed.
%    (off-peak = trivial/none is fine, as at Ch5; the CLOSED-loop value is withheld.) The "charged loop" is the
%    charge pun's home (charge = loop count, made literal — re-sound it).
%  - ±1 + SPINOR ENTIRELY OUT (Ch17 owns them). Do NOT seed spin / the fermion "two-fold" — seed only "more
%    proved residues ahead."
%  - NEVER indistinguishable ⇒ identical (box/type-identity reserved Ch24+).
%  - Charge re-sounds (proved loop count). Consensus steady (Ch6 ceiling). Continuum foundational floor reserved.
%    "verdict"/2nd-person/penalty-"fine" ABSENT.
%
% CONTINUITY: state carries from Ch15. Ch15 closed pointing here — the amplitude against a lattice "proves a loop
% residue a second time, the shadow-to-theorem crossing... made once more." Ch16 IS that second proof.
%
% PROPOSED SEAM in books/experimentation/latex/chapters/16.tex:
%   KEEP the entire locked body through line 206 ("...where three more proved loop residues make the finite model prove rather than shadow.").
%   APPEND: \bigskip / \begin{center}\rule{0.4\textwidth}{0.4pt}\end{center} / \bigskip / then the coda below. (Fresh append.)
%
% >>>>>>>>>>>>>>>>>>>>>> APPEND BEGINS (fresh \rule + coda after the body) >>>>>>>>>>>>>>>>>>>>>>

\bigskip
\begin{center}
\rule{0.4\textwidth}{0.4pt}
\end{center}
\bigskip

\noindent Some evidence a case can only name, and some it can prove --- and a proof is by far the stronger
thing. The instrument has spent most of the book naming: writing a phenomenon out in full and standing behind
the label. But a few times, only a few, it does more --- it closes a loop and shows, as settled fact, that the
loop comes back charged, a thing no arrangement of ours can talk away. Those proofs are the case's firmest
ground, and it holds only a handful of them. Here it earns another.

The amplitude that opened this part is now matter's own. What was carried as a phased tally is the very stuff of
an electron, and matter, it turns out, is a wave: sent against the regular ranks of a crystal it diffracts,
bunching into sharp peaks where a wave would and nowhere else. The count is still doing the work --- a
wavelength is only the tally of phase-turns a particle winds as it travels, and the peaks fall exactly where
those turns come home in step. The peaks stand where the extra path a wave takes to the deeper rank of atoms and
back comes to a whole number of its own wavelengths, so that every rank's reflection arrives in step and adds;
at any other angle they fall out of step and cancel. Turn the crystal and the peaks flare one after another,
each rank swinging in its turn into the closing condition --- the crystal ours to turn, the flare the world's to
answer. Fire the beam and the peaks are ours to look for; where they land, and the wave they prove, are the
world's. It is why the everyday world hides the wave and a crystal shows it: a thrown stone winds its turns so
tightly that its wavelength falls below every floor, while an electron's, at the scale of the spacing between
atoms, rises just to where a crystal can read it. And the wave the crystal proved has since become a tool of its
own: an electron's wave can be made far shorter than light's, so a beam of electrons resolves what no light
reaches, and the electron microscope --- its lens a shaped magnetic field, its illumination a matter wave ---
images atoms one at a time.

And here is the proof. Send the electron's wave once around the circuit the crystal closes --- in through the
ranks of atoms and out --- and at the right angle its phase comes home having wound a whole number of times:
the loop closes charged. This the instrument does not shadow but proves, exactly as the echo maze proved its
residue among the fields --- a finite theorem about the loop, that at the peak the loop is charged and off it
trivial, with nothing in between and no continuum called upon. The matter wave is read straight off the residue
the world will not let cancel. It is the second time in the book the instrument crosses from shadow to theorem,
the second of its small handful of proved residues; and the charge the loop carries is, once more, the loop's
own count --- charge is what the loop counts, now proved on a crystal. What that residue is worth --- which way
it is signed, what number it carries --- the book still keeps back; that a charged loop is proved is the whole
of the claim here.

So the \emph{charge} is proved again, and proved as what it always was: the count a loop carries around and
brings home. Beside the proof the chapter sets two further readings of the same matter. In one, a single
collision of light against an electron balances its books to the last unit --- all that the two carry in they
carry out, the very conservation the mechanics kept at a boundary, enforced now across one event. The same
collision run the other way --- a fast electron kicking a slow photon up to an X-ray --- lights up much of the
high-energy sky. In the other, light surrenders its energy only one whole quantum at a time, never a fraction
and never a sum pooled from several, so that below a threshold no brightness whatever frees an electron --- a
wall of rate, not of amount, named as physics with the finite model claiming only the label. Between them the
two set matter's two faces side by side: an electron, a particle, shown a wave on the crystal; light, a wave,
shown a particle in the collision --- one record read from either side, the instrument taking no side of its
own. The wave face even ran in a family: one Thomson took a prize showing the electron a particle, and his son
took another showing it a wave. And the trace is the residue that survived the loop, the peak the world would
not let cancel; whatever comes to be owed now rides on a count proved, not merely named.

A proved reading is worth more to the case than a named one --- but it is still one reading among the three. The
residue proved on a crystal is another count the three hands must in the end be brought to meet on: nearer
proof, and no nearer their meeting, for a proof binds the thing it proves, not the separate hands that must
still agree.

The instrument's handful of proofs is about to grow. The next chapter turns to the particle itself --- to what
it is that carries the wave --- and there the instrument closes not one loop but several, proving more residues
in a row than anywhere else in the book. What they are, the next chapter holds; here it is enough that more
proofs are coming. The case is still open, its charge now twice proved and still unmet by the others; the court
does not sit until the record is complete.

% <<<<<<<<<<<<<<<<<<<<<< APPEND ENDS <<<<<<<<<<<<<<<<<<<<<<
